% !PRINTED 441
\weberpage{441}
%<<CONT
verschwinden. Eine solche gerade Linie wollen wir eine \gesperrt{Hauptaxe der Substitution} $A$ nennen. Aber nur gewisse besondere Substitutionen besitzen Hauptaxen. Eine Substitution mit einer Hauptaxe hat ausser dieser noch \gesperrt{einen Pol}, der in besonderen F\"allen auch in die Hauptaxe hineinfallen kann.

Legen wir, um diese Verh\"altnisse zu \"ubersehen, die Hauptaxe in die Linie $x_1=0$, so muss, wenn $x_1=0$ ist, $y_1=0$, $y_2:y_3=x_2:x_3$ sein. Daraus ergiebt sich
\[
\begin{aligned}
y_1 &= \alpha\, x_1\\
y_2 &= \alpha'\, x_1 + \beta\, x_2\\
y_3 &= \alpha''\, x_1 + \beta\, x_3.
\end{aligned}
\]
Den ausser der Hauptaxe existirenden Pol erhalten wir, wenn wir $y_1=\alpha\,x_1$, $y_2=\alpha\,x_2$, $y_3=\alpha\,x_3$ setzen; dann ergiebt sich
\[
(\beta-\alpha)\,x_2+\alpha'\,x_1=0,\quad (\beta-\alpha)\,x_3+\alpha''\,x_1=0,
\]
und dieser Punkt wird dann und nur dann in die Linie $x_1=0$ fallen, wenn $\beta=\alpha$ ist. Dieser Fall kann bei endlichen Gruppen nicht eintreten (ausser bei der Multiplication).

Wenn ein von der Hauptaxe verschiedener Pol existirt, so k\"onnen wir diesen in die Ecke $x_2=0$, $x_3=0$ legen, und die Substitution erh\"alt die Form
\[
y_1=\alpha\,x_1,\quad y_2=\beta\,x_2,\quad y_3=\beta\,x_3,
\]
wo $\alpha$ von $\beta$ verschieden ist.

Eine gerade Linie $u_1x_1+u_2x_2+u_3x_3=0$ ist dann und nur dann ihr eigenes Bild, wenn entweder $u_2=u_3=0$ oder $u_1=0$ ist.

\gesperrt{Wenn also eine Substitution einen Pol und eine Hauptaxe hat, so bleiben ausser dieser alle geraden Linien und nur die unge\"andert, die durch den Pol gehen.}

Es ist noch zu bemerken, dass, wenn drei getrennte Pole vorhanden sind, diese nicht in eine gerade Linie fallen k\"onnen, ohne dass die Linie zur Hauptaxe wird. Denn wenn eine gerade Linie ihr eigenes Bild ist, so k\"onnen, wenn sich die Linie nicht Punkt f\"ur Punkt selbst entspricht, nur zwei Punkte auf ihr liegen, die ihr eigenes Bild sind.

Fassen wir dies Alles zusammen, so finden wir folgende Arten von tern\"aren linearen Substitutionen, wobei zusammenfallende Pole nur einmal mitgez\"ahlt sind:
%CONT>>
